\cotefile{162-5}
\scanpage{31}
les deux notions de connexion se correspondent biunivoquement, de $=$ que les deux notions de \inserted{tenseur courbure, de} connexion intégrable \ldots, et bien sur les complexes
\[
\bigwedge^*\Omega\otimes M\qquad\text{et}\qquad C^*(\mathscr T,M).
\]

11.\quad Sous les conditions du th de Quillen, ($X_0$ de car. nulle, \deleted{\unsure{la}} $(\bigwedge^*\Omega,\delta)$ \deleted{\unsure{loc. libre}}\inserted{avec $\Omega$ plat $\Leftarrow$} correspondant à une catégorie formelle $\mathcal C$ avec $\Omega$ plat, pour un Module $M$ sur $\underline O_{X_0}$, les structures de $C$-descente sur $M$ correspondent aux \deleted{\unsure{struct.}} $\Omega$-connexions sur $M$.

\margin{généralise mon th. fondamental de démarrage des cristaux\,\ldots. \unsure{Peut-il|Peut il} s'y ramener, en raisonnant dans un topos convenable ?!}

De plus, la cohomologie de De Rham de $M$, i.e $\mathbb H^{\unsure{\alpha|*}}(X,\bigwedge^{\unsure{\alpha|2|*}}\Omega\otimes M)$, est can. isom à la coh. de Hochschild (ou Čech-Alexander) de $C$ à coeff.\unsure{,|.|-} \unsure{dans|\deleted{dans}} $M$, plus précisément il y a un quasi-isom \inserted{canonique} de \inserted{complexes de} faisceaux abéliens (\unsure{dans 2 catég.|les 2 catég.} \unsure{dérivées|dérivée})
\[
\bigwedge^*\Omega\otimes M\simeq C^*(\mathcal C,M)
\]

\uline{Variante à puissances divisées}.

\margin{30}
